\documentclass[11pt]{article}

\usepackage[margin=1in]{geometry}
\usepackage{amsmath,amssymb,booktabs,array}
\usepackage{microtype}
\microtypesetup{expansion=false}
\usepackage{xcolor}
\usepackage{listings}
\usepackage{hyperref}
\usepackage{enumitem}
\usepackage{fancyhdr}

\definecolor{navy}{HTML}{264653}
\definecolor{teal}{HTML}{2A9D8F}
\definecolor{paper}{HTML}{F7F7F5}
\hypersetup{
  colorlinks=true,
  linkcolor=navy,
  citecolor=teal,
  urlcolor=teal,
  pdftitle={PyBED Technical Report}
}
\lstset{
  language=Python,
  basicstyle=\ttfamily\small,
  keywordstyle=\color{navy}\bfseries,
  commentstyle=\color{teal},
  backgroundcolor=\color{paper},
  frame=single,
  rulecolor=\color{black!15},
  columns=fullflexible,
  keepspaces=true,
  showstringspaces=false,
  breaklines=true,
  xleftmargin=0.5em,
  framexleftmargin=0.4em
}
\setlist{nosep,leftmargin=1.5em}
\setlist[itemize]{label=$\bullet$}
\pagestyle{fancy}
\setlength{\headheight}{14pt}
\fancyhf{}
\lhead{PyBED v0.2}
\rhead{Technical report}
\cfoot{\thepage}

\title{\textbf{PyBED: A Small Interface for Bayesian Experimental Design and Inverse Problems}}
\author{PyBED contributors}
\date{Version 0.2.0 -- September 2026}

\begin{document}
\maketitle

\begin{abstract}
Bayesian experimental design methods often use the same scientific problem but
reimplement its simulator, histories, random draws, metrics, and saved data. Those
differences make method comparisons hard to read and harder to trust. PyBED keeps
this shared layer small. An environment combines a prior, a simulator, shape
information, and optional design, inference, prediction, likelihood, candidate-set,
and joint-model callables. A batch stores complete observations
$o_t=(x_t,y_t)$, optional particle beliefs, targets, masks, and problem context.
The same interface covers likelihood-based DAD, likelihood-free iDAD, joint models
such as ALINE and JADAI, task-driven Action-BED, PSST-style particle transport,
and amortized Bayesian inverse problems trained from joint samples. This report
defines the implementation as shipped in PyBED 0.2, documents every benchmark and
data path, and explains the included plain-PyTorch and Lightning/W\&B examples.
\end{abstract}

\tableofcontents
\newpage

\section{Problem and scope}

A sequential experiment has an unknown quantity $\theta$, a design $x_t$, an
outcome $y_t$, and the complete observation
\begin{equation}
  o_t=(x_t,y_t).
\end{equation}
The history after $t$ steps is $h_t=(o_1,\ldots,o_t)$. A design policy uses
$h_t$ to choose $x_{t+1}$. An inference model uses the same history to return a
posterior distribution, a point, logits, or particles. A prediction model may
instead target a future outcome or another downstream quantity.

An ordinary inverse problem is the no-design special case
\begin{equation}
  y=G(\theta,\eta),
\end{equation}
where $x$ is absent and $o=(\varnothing,y)$. Keeping that case in the same data
model matters: simulation-based inference and BED differ in when a design is
chosen, not in the basic meaning of paired parameter and observation data.

PyBED standardizes the parts that can reasonably be shared:
\begin{itemize}
  \item prior sampling and optional likelihood evaluation;
  \item batched simulation, histories, masks, and targets;
  \item adaptive and non-adaptive data collection;
  \item particle beliefs, candidate sets, and policy sample information;
  \item reproducible saved datasets, replay, scores, figures, and run records.
\end{itemize}
It does not prescribe a neural architecture, optimizer, information estimator,
reinforcement-learning algorithm, neural operator, or training service.

\section{The public data model}

\subsection{Observation names}

The main storage names are \texttt{x} and \texttt{y}. Familiar longer names remain
available, while \texttt{o} and \texttt{obs} mean the complete pair:
\begin{lstlisting}
batch = pb.Batch(theta=theta, x=x, y=y)

batch.x          # design
batch.design     # the same design
batch.y          # measured outcome
batch.outcome    # the same outcome
batch.o          # Observation(x, y)
batch.obs        # the same complete observation
\end{lstlisting}
An \texttt{Observation} is a small named pair. It can be unpacked as
\texttt{x, y = batch.o}; unlike concatenation, it remains valid when $x$ is a
vector and $y$ is an image. For inverse problems, \texttt{batch.x is None} and
the outcome remains in \texttt{batch.y}.

This is a deliberate change from version 0.1, where \texttt{Batch.obs} meant only
$y$. Saved version-1 datasets remain readable; newly saved datasets use explicit
\texttt{x} and \texttt{y} fields.

\subsection{Batch fields and shapes}

\begin{center}
\begin{tabular}{>{\ttfamily}p{0.16\linewidth}p{0.72\linewidth}}
\toprule
Field & Meaning \\
\midrule
theta & Latent truth, optional when real observations have no known truth. \\
x & Designs with usual shape $(B,T,\ldots)$; optional for inverse problems. \\
y & Outcomes with usual shape $(B,T,\ldots)$, or $(B,\ldots)$ without a design sequence. \\
mask & Boolean valid-step mask of shape $(B,T)$. \\
target & Label, selected parameter, image, predictive value, or downstream action target. \\
belief & Optional batched \texttt{ParticleCloud}. \\
context & Per-example masks, dimensions, physical settings, or other model inputs. \\
meta & Dataset-level provenance such as the seed, budget, and epoch number. \\
\bottomrule
\end{tabular}
\end{center}

\texttt{Batch.history(steps)} slices axis 1 and keeps all later event axes. This
works for scalar designs, sets of sensors, patches, and other structured events.
\texttt{Batch.index}, \texttt{Batch.to}, and \texttt{Batch.numpy} also walk
nested context and particle fields. Per-example context is therefore indexed and
collated with the data rather than silently remaining at its original batch size.

\section{Environment and component calls}

\subsection{The BED value}

A \texttt{BED} object holds a name, prior, simulator, event specifications, budget,
and optional components. \texttt{Spec(shape)} describes trailing event axes only:
\texttt{Spec((2,))} accepts one two-dimensional design or any leading batch and
time axes followed by two.

The four single-purpose verbs are:
\begin{lstlisting}
y = env.simulate(theta, x)
x = env.design(history)
inference = env.infer(history)
prediction = env.predict(query, history)
\end{lstlisting}
\texttt{simulate} also accepts no $x$. \texttt{infer} intentionally returns
\texttt{Any}: a point estimator is no less valid than a distribution-valued
posterior, and an empirical cloud is not forced into a density interface.

Components are ordinary callables:
\begin{lstlisting}
method = env.with_components(
    simulate=my_simulator,
    design=my_policy,
    infer=my_posterior,
    predict=my_predictor,
    log_prob=my_likelihood,
    candidates=my_candidates,
)
\end{lstlisting}
The original environment is unchanged. PyBED inspects callable signatures and only
passes named context that the callable accepts. There is no base model class.

\subsection{One-pass joint models}

ALINE and the shared PSST policy can update inference and choose a design in one
backbone pass. Calling separate heads through separate public verbs would repeat
that work. PyBED therefore provides three combined calls:
\begin{align*}
 \texttt{infer\_design}(h)
   &\longrightarrow (\text{inference},x),\\
 \texttt{design\_predict}(h)
   &\longrightarrow (x,\text{prediction}),\\
 \texttt{design\_infer\_predict}(h)
   &\longrightarrow (x,\text{inference},\text{prediction}).
\end{align*}
The word order is the output order. A model attaches only the calls it computes
together:
\begin{lstlisting}
method = env.with_components(
    infer_design=model.infer_design,
    design_infer_predict=model.design_infer_predict,
)
\end{lstlisting}
An attached 	exttt{design_infer_predict} component also serves either two-output
verb by returning the requested members. Otherwise, the environment composes the
available single-purpose calls. Thus separate models and genuinely shared models
use one public surface without pretending their compute graphs are the same.

When \texttt{Stream} uses an attached joint model, it calls that model once at each
adaptive step. A returned \texttt{ParticleCloud} becomes the next history's belief.

\subsection{Likelihoods and score-function policies}

An optional observation likelihood has the direct order
\begin{lstlisting}
log_p = env.log_prob(y, theta, x)
\end{lstlisting}
and may omit $x$ for an inverse problem. The location and MNIST-patch environments
provide their Gaussian likelihoods. Likelihood-free environments simply leave this
component absent.

A deterministic policy returns $x$. A sampled policy may return
\begin{lstlisting}
pb.PolicySample(x=x, log_prob=log_prob, entropy=entropy)
\end{lstlisting}
so REINFORCE-like training retains the terms produced by the actual sampling step.
PyBED does not define the reward, baseline, discounting, or credit assignment,
because those choices belong to the method.

\subsection{Candidate sets}

\texttt{candidate\_pool} may be a fixed tensor, a batched tensor, or a callable of
the current history. \texttt{env.candidates(history)} returns it, and PyBED passes
it to policies and joint calls under the clear argument name \texttt{candidates}.
A random policy samples uniformly from a supplied pool and returns its log
probability. Continuous policies are unaffected.

\section{Particle methods}

\texttt{ParticleCloud} keeps
\begin{equation}
  C=(\{\theta^{(s)}\}_{s=1}^{S},\{w_s\}_{s=1}^{S},m_\theta),
\end{equation}
where weights and a parameter mask are optional. The class computes a weighted
mean, effective sample size, and weighted resampling, and supports device and NumPy
conversion. Batched storage conventionally has shape
$(B,S,\ldots)$ and \texttt{particle\_dim=1}; an unbatched empirical prior uses
\texttt{particle\_dim=0}.

Two roles are kept separate:
\begin{itemize}
  \item \texttt{EmpiricalPrior} resamples a fixed initial cloud to draw truths;
  \item \texttt{Batch.belief} carries a changing, per-episode belief used by a
  model during sequential inference and design.
\end{itemize}
This distinction is necessary for PSST. Its current cloud $C_t$ is a model state,
not merely the probability law that generated the hidden truth once.

For a categorical target, particles may be labels or one-hot vectors. If
$z^{(s)}\in\{e_0,\ldots,e_9\}$, then
\begin{equation}
  \widehat p(z=k\mid h)=\sum_s \bar w_s z^{(s)}_k.
\end{equation}
\texttt{metrics.class\_probabilities} returns these normalized counts. PyBED does
not sharpen or calibrate them.

\section{Data generation and reuse}

\subsection{Seeds}

\texttt{Seeds(seed)} creates a repeatable Torch or NumPy generator from a short
label and integer indices. Separate labels are used for truths, designs, simulator
noise, and loader order. This makes a small code change in one part of an
experiment less likely to change every later random draw.

For a fixed environment, root seed, and epoch, one \texttt{Stream} caches the same
ordered latent truths for each policy. \texttt{exp.compare} checks that equality
before reporting policy scores.

\subsection{Stream}

\texttt{Stream(env, episodes, seed, budget)} materializes a finite dataset.
Adaptive policies are called step by step because $x_{t+1}$ depends on $h_t$.
For a non-adaptive experiment, \texttt{Stream} instead:
\begin{enumerate}
  \item creates the whole tensor $x_{1:T}$, or accepts it from the user;
  \item adds the time axis to $\theta$ and per-example context;
  \item calls a vectorized simulator once.
\end{enumerate}
This is used automatically when there is no policy. Custom simulators that cannot
accept a leading time axis set \texttt{vectorized=False} and use the clear
step-by-step path.

\begin{lstlisting}
stream = pb.data.Stream(env, episodes=8192, seed=42)
random_data = stream.generate()
fixed_data = stream.generate(x=fixed_x)
adaptive_data = stream.generate(policy)
\end{lstlisting}

\subsection{Dataset files and replay}

\texttt{EpisodeDataset} requires paired $\theta$ and $y$ but no longer requires
$x$. This directly represents the i.i.d. joint data used by amortized inverse
methods. Version-2 files store $\theta$, $x$, $y$, masks, targets, beliefs, nested
context, metadata, and a readable JSON shape manifest. Version-1 files are upgraded
while loading.

\texttt{ReplayBuffer(capacity)} keeps recent detached Batch rows. Adding a large
batch splits it into rows; sampling collates rows back into a batch. This supports
PSST-style reuse of truth, evidence, predicted particles, shape masks, and policy
terms without defining a PSST-specific buffer. It complements rather than replaces
a stable saved dataset.

\section{Reference environments}

\subsection{Location finding, including mixed shapes}

For source $k$ at $\theta_k\in\mathbb R^d$ and sensor $x\in\mathbb R^d$,
\begin{equation}
  \mu(\theta,x)
  =b+\sum_{k=1}^{K}\frac{\alpha_k}
  {m+\lVert\theta_k-x\rVert_2^2},
  \qquad
  y=\log\mu(\theta,x)+\sigma\epsilon.
\end{equation}
All constants and the log transform are configurable. The implementation is
differentiable in $x$ and $\theta$ and provides the Gaussian log likelihood used
by DAD.

Fixed problems use integer \texttt{sources} and \texttt{dims}. A single training
run may instead use choices:
\begin{lstlisting}
env = pb.make(
    "location-v0",
    sources=(1, 2, 4),
    dims=(1, 2, 3),
)
\end{lstlisting}
The prior samples a source count and dimension for each episode, pads to the largest
shape, and records \texttt{num\_sources}, \texttt{source\_dim},
\texttt{source\_mask}, and \texttt{theta\_mask}. Masked sources and coordinates
make no contribution to the signal. A model may use the same masks for attention.

\subsection{MNIST discovery and classification}

\texttt{mnist-discovery-v0} returns a smooth masked full image at each continuous
centre. It is useful for posterior image reconstruction.

\texttt{mnist-classification-v0} implements the Action-BED task
\cite{rossa2026actionbed}. The latent parameter is an image
$\theta\in[0,1]^{1\times28\times28}$ and the target is its label
$z\in\{0,\ldots,9\}$. A design $x\in[0,1]^2$ gives a continuous top-left
patch location. Bilinear grid sampling returns
\begin{equation}
  y=A_x\theta+\eta,\qquad
  \eta\sim\mathcal N(0,\sigma^2 I_{25}).
\end{equation}
Zero padding applies outside the image and the default horizon is five. Bilinear
sampling keeps the simulator differentiable in $x$. The image and label are sampled
together, so a batch cannot accidentally mismatch them.

This is a hierarchical problem. The available likelihood is
$p(y\mid\theta,x)$, while the class likelihood
\begin{equation}
  p(y\mid z,x)=\int p(y\mid\theta,x)p(\theta\mid z)\,d\theta
\end{equation}
is implicit under the empirical image bank. DAD can target image information with
the explicit likelihood; Action-BED, ALINE, and classification critics can train
directly for $z$.

\subsection{Fields, dynamics, and preferences}

\texttt{advdiff-v0} infers an initial field from time series at designed sensors.
Its simulator uses semi-Lagrangian advection, a five-point explicit diffusion step,
an optional reaction, bilinear sensor sampling, and Gaussian noise. It checks the
explicit diffusion stability ratio. It is a compact differentiable reference
solver, not a paper-exact finite-element mesh.

\texttt{pendulum-v0} infers gravity and damping from a stochastic
Euler--Maruyama trajectory whose initial angle and velocity are designed.
\texttt{ces-v0} observes noisy preferences between designed baskets.
\texttt{death-v0} observes the affected count of a pure-death process at a designed
time. These broaden simulator shapes and differentiability conditions without
adding method-specific code.

\section{Generic amortized inverse problems}

\texttt{envs.inverse} accepts a prior, a forward simulator, and the shapes of
$\theta$ and $y$:
\begin{lstlisting}
env = pb.envs.inverse(
    prior=prior,
    simulator=forward,
    theta_shape=(64, 64),
    y_shape=(20,),
)
pairs = pb.data.Stream(env, 8192, seed=4).generate()
\end{lstlisting}
External paired data are equally direct:
\begin{lstlisting}
pairs = pb.data.EpisodeDataset(pb.Batch(theta=u, y=o))
\end{lstlisting}

Baptista, Kaveh, and Stuart assume $N$ i.i.d. samples
\begin{equation}
  \{(\theta^{(i)},y^{(i)})\}_{i=1}^{N}\sim\gamma,
\end{equation}
where $\gamma$ is the joint law induced by the prior and forward model
\cite{baptista2026energy}. Their observation-dependent map
$T_\phi(z;y)$ pushes a reference law toward $p(\theta\mid y)$. Its average
energy-score objective is
\begin{align}
 \mathcal J(\phi)
 ={}&\mathbb E\lVert T_\phi(z;y)-\theta\rVert \\
 &-\frac12\mathbb E\lVert
 T_\phi(z;y)-T_\phi(z';y)\rVert .
\end{align}
Only joint samples and reference samples are needed; no likelihood is evaluated.

\texttt{metrics.energy\_score} implements this differentiably for tensor particles
and weighted \texttt{ParticleCloud} objects. The option
\texttt{unbiased=True} removes self-pairs from the repulsion term. If the prior is
not separately available, distinct parameter rows from the joint sample can serve
as empirical reference draws; method code is responsible for the required distinct
index construction.

Unknown functions are represented by their discretized tensor axes. A
Cameron--Martin-informed map such as
\begin{equation}
  T_\phi(u;y)=u+C^{1/2}S_\phi(u;y)
\end{equation}
belongs in the model: $C^{1/2}$ and the neural operator depend on the scientific
prior, basis, and grid. PyBED supplies paired data, fields, loaders, particles, and
the score without claiming that a finite tensor interface alone proves a
function-space measure result.

The executable \texttt{examples/inverse/energy\_transport.py} demonstrates a
likelihood-free $y=\theta^2+\eta$ problem with a residual transport.

\section{How the method families fit}

\subsection{DAD}

DAD samples a generating $\theta_0$ and prior contrastives
$\theta_1,\ldots,\theta_L$, rolls out one adaptive history, and accumulates the
same observed-history likelihood under every candidate. PyBED supplies
\texttt{env.log\_prob}; the objective is
\begin{equation}
 \widehat I_{\mathrm{sPCE}}
 =\frac1B\sum_b\left[
 \ell_b(\theta_0)-\log\left(
 \frac1{L+1}\sum_{j=0}^{L}e^{\ell_b(\theta_j)}
 \right)\right].
\end{equation}
\texttt{metrics.spce} consumes the $(B,L+1)$ likelihood table. The included
\texttt{dad\_lightning\_wandb.py} example differentiates through the location
simulator and maximizes this bound.

\subsection{iDAD}

iDAD uses the same $o_{1:T}$ batches and policy calls, but replaces explicit
likelihoods with a critic and an InfoNCE or NWJ objective
\cite{ivanova2021idad}. The critic and negative-sample construction remain in the
method. Nothing in \texttt{Batch} requires likelihood access.

\subsection{ALINE}

ALINE can attach \texttt{infer\_design} or
\texttt{design\_infer\_predict}. A candidate provider supplies its allowed query
set; a categorical policy returns \texttt{PolicySample}; the inference member may
be a GMM, a point, a predictive distribution, or class logits. \texttt{target}
holds the requested parameter subset or predictive task, while attention details
stay in the ALINE model \cite{huang2025aline}.

\subsection{JADAI}

JADAI attaches a posterior and differentiable design policy, separately or through
a combined call \cite{bracher2025jadai}. Images and fields remain structured event
axes. Reparameterized built-in simulators permit gradients from posterior quality
through $y$ and $x$.

\subsection{Action-BED}

Action-BED optimizes a downstream Bayes action rather than requiring a posterior
density \cite{rossa2026actionbed}. The downstream truth sits in
\texttt{Batch.target}. The location example jointly trains a design policy and
source estimator. The MNIST environment supplies the classification target, for
which a downstream action is a ten-dimensional logit or probability vector.

\subsection{PSST}

PSST places its current particles in \texttt{Batch.belief}. Its shared update and
proposal model attaches \texttt{infer\_design}; the output is the new cloud and
next $x$. Padded location masks allow varying source counts and dimensions.
\texttt{ReplayBuffer} can retain recent clouds with their truth and observations.
\texttt{energy\_score} accepts the physical posterior particles directly. This
supports both one-shot prior-to-posterior transport and repeated Bayes-map use
without calling either one a conventional density.

\section{Metrics}

The metric module contains MSE, RMSE, log-MSE, negative log score, interval
coverage, rank calibration error, exact small-source set MSE, Monte Carlo expected
information gain, sPCE, cross entropy, and classification accuracy.

For weighted particles $z_s$ with normalized weights $w_s$, the energy score is
\begin{equation}
  \operatorname{ES}
  =\sum_s w_s\lVert z_s-\theta^*\rVert^\beta
  -\frac12\sum_{s,r}w_sw_r\lVert z_s-z_r\rVert^\beta .
\end{equation}
The unbiased option renormalizes off-diagonal pair weights. Parameter masks are
applied before distances. \texttt{posterior\_summary} collects common scalar
diagnostics, while classification metrics accept logits or categorical particles.

Metrics do not automatically turn a point into a distribution or sharpen a cloud.
That keeps representation choices visible in reported comparisons.

\section{Experiments, Lightning, and W\&B}

\texttt{exp.Run} writes a plain local directory containing normalized
configuration, system and Git metadata, JSONL metrics, checkpoints, figures,
datasets, and a final summary. Checkpoints include model, optimizer, and Python,
NumPy, CPU-Torch, and CUDA random states. \texttt{Run.log} returns the row it
writes, so plain PyTorch code may also call
\texttt{wandb\_run.log(row)} without an adapter.

Lightning and W\&B are an optional extra rather than base dependencies:
\begin{lstlisting}
pip install -e '.[lightning]'

WANDB_MODE=offline python examples/dad/dad_lightning_wandb.py

PYBED_DEVICES=4 WANDB_MODE=online python examples/dad/dad_lightning_wandb.py
\end{lstlisting}
The example uses a \texttt{LightningModule}, automatic accelerator selection,
DDP when more than one GPU is requested, Lightning checkpoints, and
\texttt{WandbLogger}. It logs training sPCE and an evaluation trajectory. Setting
\texttt{WANDB\_MODE=disabled} gives a fast service-free smoke run.

This split keeps two useful paths:
\begin{itemize}
  \item small studies may use the readable local \texttt{exp.train} loop;
  \item multi-GPU studies may use Lightning without a PyBED wrapper around its
  trainer, callbacks, or logger.
\end{itemize}

\section{Backends and module layout}

Built-in simulators use PyTorch. \texttt{backends.convert} explicitly converts
arrays, batches, nested context, policy samples, and particle clouds to Torch,
NumPy, or optional JAX. It uses DLPack where practical. PyBED does not silently
translate numerical solvers across frameworks, because the result may not have the
same numerical behavior.

\begin{center}
\begin{tabular}{lp{0.70\linewidth}}
\toprule
Module & Contents \\
\midrule
\texttt{core} & BED, Batch, Observation, ParticleCloud, PolicySample, specs, priors, registry. \\
\texttt{sim} & Batched location, image, PDE, SDE, preference, and event simulators. \\
\texttt{envs} & Benchmark factories and the generic inverse-problem factory. \\
\texttt{data} & Stream, Seeds, datasets, loaders, serialization, replay. \\
\texttt{metrics} & Posterior, information, calibration, set, and classification scores. \\
\texttt{vs} & Prior, trajectory, field, image, training, and benchmark plots. \\
\texttt{exp} & Local runs, checkpoints, a small training loop, policy comparison. \\
\texttt{backends} & Explicit array and container conversion. \\
\bottomrule
\end{tabular}
\end{center}

\section{Validation and limits}

The tests cover shape checks, aliases, each joint-call order, component replacement,
particle statistics, weighted and unbiased energy scores, candidate sampling,
variable location masks, nested context collation, replay, vectorized one-call
simulation, fixed designs, inverse pairs without $x$, MNIST patches and labels,
simulator gradients, likelihoods, serialization, policy comparison, checkpoints,
and quick end-to-end examples. The Action-BED, energy-transport, and optional
Lightning/W\&B scripts each have a short configuration selected by
\texttt{PYBED\_QUICK=1}.

Important limits remain:
\begin{itemize}
  \item a compact built-in environment is not automatically a paper-exact
  reproduction;
  \item distributed on-the-fly simulator workloads may need method-specific data
  placement even though distributed model training works through Lightning;
  \item high-source-count set matching needs an assignment solver rather than the
  included exact small-set permutation metric;
  \item a tensor discretization does not by itself preserve infinite-dimensional
  absolute continuity;
  \item fair comparisons must still report model size, optimizer, training budget,
  contrastive count, particle count, seeds, and estimator uncertainty.
\end{itemize}

\section{Conclusion}

PyBED 0.2 treats an observation as the pair $(x,y)$, treats particles as a
first-class inference representation, and treats no-design inverse problems as a
normal part of the same interface. Joint verbs avoid repeated backbone work for
ALINE- and PSST-like models. Candidate pools, score-function values, changing
parameter masks, replay, and vectorized non-adaptive simulation cover the practical
gaps exposed by those methods. Optional Lightning and W\&B examples provide
multi-GPU training and tracking without moving either dependency into the core.
The result remains a small research library whose main objects can be read without
learning a software framework first.

\begin{thebibliography}{9}

\bibitem{foster2021dad}
A. Foster, D. R. Ivanova, I. Malik, and T. Rainforth.
Deep Adaptive Design: Amortizing Sequential Bayesian Experimental Design.
\emph{Proceedings of ICML}, 2021.
\url{https://proceedings.mlr.press/v139/foster21a.html}.

\bibitem{ivanova2021idad}
D. R. Ivanova, A. Foster, S. Kleinegesse, M. U. Gutmann, and T. Rainforth.
Implicit Deep Adaptive Design: Policy-Based Experimental Design without Likelihoods.
\emph{Advances in Neural Information Processing Systems}, 2021.

\bibitem{huang2025aline}
D. Huang et al.
ALINE: Joint Amortization for Bayesian Inference and Active Data Acquisition.
arXiv:2506.07259, 2025.

\bibitem{bracher2025jadai}
N. Bracher et al.
Jointly Amortizing Adaptive Design and Bayesian Inference.
arXiv:2512.22999, 2025.

\bibitem{rossa2026actionbed}
L. Rossa, A. Phillips, and T. Rainforth.
Task-Driven Bayesian Experimental Design with Singly Intractable Objectives.
2026.

\bibitem{baptista2026energy}
R. Baptista, H. Kaveh, and A. M. Stuart.
Energy-Based Transport for Amortized Bayesian Inference.
arXiv:2605.15407, 2026.

\end{thebibliography}

\end{document}
